% Plan v4 / section B: the anchor-scaling law. Written 2026-09-07 on branch iclr-2027.
% Every number here comes from results/anchor_scaling_summary.csv, _paired.csv and
% budget_drift.csv, produced by analysis/anchor_scaling.py and analysis/budget_drift.py.

\section{The separation widens as the safe model improves}\label{sec:scaling}

Proposition~\ref{prop:threshold} makes the certificate's fate turn on one comparison: the budget
against the safe model's surprisal of the protected work. Whether a deployer can pick a budget that
is both informative and generous therefore depends on how $s(x)$ compares with $c_{\mathrm{use}}$,
the rate ordinary traffic diverges from the same safe model. Both are properties of the models and
the content, settled before any mechanism is designed. Neither has been measured across safe models,
and the question has an obvious practical edge: safe models are getting better, and it is not
obvious whether that helps.

\paragraph{Setup}
We score ten safe models spanning $0.17$B to $7$B parameters and three independently curated,
openly licensed corpora: KL3M (US public-domain legal text), Pleias~1.0 (Common Corpus) and Comma
(the Common Pile), the last including the $1$T- and $2$T-token checkpoints of one $7$B model, which
vary training data at fixed capacity. Every model scores the same $758$ protected passages, the
same $50$ public-domain control books, and the same $300$ generations the risky model produced on
ordinary instruction prompts. Rates are nats per character throughout, because total string
log-probability is tokenizer-invariant and these models carry $64$k, $128$k and $152$k
vocabularies; that is also why the measurement needs no shared vocabulary, and so is not restricted
to safe models that could actually be fused with the risky one.

\paragraph{Result}
Both rates fall as the safe model improves, and the one that must be permitted falls faster than
the one that must be forbidden. Across the widest span in each family:
\begin{center}\small
\begin{tabular}{@{}llrrrr@{}}
\toprule
corpus & span & $\Delta c_{\mathrm{use}}$ & $\Delta s(x)$ & margin & sign test \\
\midrule
Common Corpus & $0.35 \to 3$B   & $-25.8\%$ & $-9.8\%$  & $3.45 \to 4.19$ & $16/16$, $p = 3.1\times10^{-5}$ \\
KL3M          & $0.17 \to 3.7$B & $-31.3\%$ & $-16.0\%$ & $1.94 \to 2.37$ & $16/16$, $p = 3.1\times10^{-5}$ \\
Common Pile   & $1.8 \to 7$B    & $-28.6\%$ & $-12.0\%$ & $4.07 \to 5.02$ & $16/16$, $p = 3.1\times10^{-5}$ \\
\bottomrule
\end{tabular}
\end{center}
The margin $s(x)/c_{\mathrm{use}}$ rises monotonically with capability in every family
(Figure~\ref{fig:frontier}). Because every model scores the same novels and the same generations
the comparison is paired, and the margin rises for all sixteen novels in all three families; it
also rises with training tokens at fixed capacity, $4.63$ to $5.02$ from the $1$T to the $2$T
checkpoint of the same $7$B model. The mechanism is unsurprising once stated: a better safe model
closes the gap to the risky model faster on generic instruction-following output than on a
particular novel, because the novel is idiosyncratic and the assistant register is not.

This corrects the earlier audit~\citep{vijayavallabh2026audit}, which observed that a $7$B Common
Pile anchor assigns the protected passages fewer nats than a $1.8$B one and concluded that scaling
the anchor weakens the guarantee. The observation is right, $0.778 \to 0.685$ nats per character;
the conclusion does not follow, because $c_{\mathrm{use}}$ falls further over the same span,
$0.191 \to 0.137$. Of the ten safe models scored here the $7$B anchor has the largest margin, not
the smallest. Both readings of the measurement are true and belong together: in absolute terms the
range of budgets over which the certificate says anything shrinks as anchors improve, so a deployer
who fixed $k$ in nats will find it going vacuous under a better anchor; in relative terms the two
rates separate. Which matters depends on whether the budget is chosen in nats or as a fraction of
what ordinary traffic costs, and only the second is available to a deployer who does not know the
protected works in advance.

The trend has a terminus: it is driven by $c_{\rm use}$ falling, so in the limit the decoder becomes
the identity on ordinary prompts and the mechanism's value vanishes in exactly the regime that makes
its guarantee comfortable.

\paragraph{Two confounds, both measured}
A safe model could look good on protected text merely by being bad at prose, and every safe model
here is a base model while the risky model is instruction-tuned. Both are controlled in
Appendix~\ref{app:scaling}: the ratio of public-domain control to protected surprisal is $0.630$ to
$0.887$ for every model, so the separation is not general fluency, and a $1$B \emph{base} risky
model roughly halves $c_{\rm use}$ and raises every margin, making the numbers above the
conservative end. In all nine corpus-by-risky-model combinations the margin rises $16$ of $16$.

\paragraph{What the margin does not license}\label{sec:price}
The margin says how far apart the two rates are; it does not say a budget between them leaves
utility intact, because the rate a decoder actually pays is not $c_{\mathrm{use}}$ but the
$5$-to-$32$-times smaller rate reported in §\ref{sec:frontier}, measured across six prompt classes
and two budgets. Cheap-step selection does not account for that gap, since the decoder pays more
per unaltered step than the cheapest matching fraction of the unconstrained distribution would. The
budgeted decoder walks a different trajectory from the one the risky model would have walked, and
its price is a property of that trajectory --- the approximation gap Theorem~\ref{thm:nfl} leaves
open, and the reason no quantity computable before serving predicts it.
